\chapter{Methodology}

\section{Overview}

The methodology is organized around three main research components that correspond to the identified research gap:

\begin{enumerate}
\item construction of a LiDAR drone data and benchmark framework,
\item development of a tracking and state-estimation pipeline from sparse LiDAR detections,
\item development and evaluation of short-horizon trajectory prediction methods.
\end{enumerate}

These components form a pipeline. The first component provides realistic and controllable data. The second component transforms per-frame detections into temporally consistent drone tracks. The third component predicts future drone motion from these tracks. The overall goal is to evaluate how sparse LiDAR observations can be transformed into safety-relevant state and trajectory information for helicopter detect-and-avoid perception.

\section{Component 1: LiDAR Drone Data and Benchmark Framework}

\subsection{Purpose}

The first methodological component establishes the data basis for the dissertation. Since public LiDAR datasets for small airborne drone encounters are limited, a controlled data generation and benchmarking framework is required. The purpose of this component is to generate datasets that contain realistic sparse point-cloud observations, accurate ground-truth object states, and temporally consistent instance annotations.

The benchmark must support three tasks simultaneously: object detection, multi-object tracking, and trajectory prediction. Therefore, the data must include not only point clouds and 3D bounding boxes, but also object identities, timestamps, scene structure, and temporally linked annotations.

\subsection{Synthetic LiDAR Data Generation}

A simulation environment will be used to generate LiDAR observations of drone encounters. The sensor model will include parameters such as maximum range, vertical field of view, angular resolution, range noise, dropout, and intensity variation. The purpose is not to create visually realistic scenes alone, but to create point-cloud statistics that are useful for studying sparse airborne object perception.

The generated scenarios will vary the number of drones, distance to the sensor, relative altitude, crossing angle, approach speed, and trajectory type. Special attention will be given to altitude variation because flying drones are not constrained to the ground plane. This is important because many efficient 3D detectors use bird's-eye-view or pillar-based representations, which can be strong in horizontal localization but weaker in vertical localization \cite{lang2019pointpillars,shi2022pillarnet}. The dataset will therefore deliberately include scenarios where drones occupy similar horizontal positions but different heights.

The data generation process will also include negative and hard-negative examples. Empty scenes, background clutter, distant drones with few returns, and partial visibility cases are necessary to evaluate false positives and robustness. This is important because a detect-and-avoid perception system must be reliable not only in ideal target-present cases, but also under uncertainty and absence of targets.

\subsection{Dataset Format and Splitting}

The generated data will be organized in a nuScenes-like format because this structure is supported by common LiDAR detection toolchains and allows temporal linking of samples, sample data, ego poses, and annotations \cite{caesar2020nuscenes,openpcdet2020}. Each frame will contain a LiDAR point cloud, sensor pose, 3D object annotations, instance identifiers, and timestamps. Annotation chains will be preserved so that drone trajectories can be reconstructed from the dataset.

The dataset will be split at the level of runs or scenarios rather than individual frames. This prevents temporally adjacent frames from appearing in both training and test data. The intended split consists of:

\begin{itemize}
\item a training set for model fitting,
\item a validation set for parameter selection,
\item an unseen benchmark set for final performance assessment.
\end{itemize}

This scenario-level separation is essential because trajectory prediction and tracking must be evaluated on unseen motion patterns, not merely on neighboring frames from the same sequence.

\subsection{Baseline Detection Models}

The first experimental baseline will use established LiDAR 3D object detection architectures. Pillar-based models such as PointPillars and PillarNet are particularly relevant because of their computational efficiency and suitability for real-time applications \cite{lang2019pointpillars,shi2022pillarnet}. Voxel- and point-based alternatives such as SECOND, PV-RCNN, and CenterPoint will be considered as comparison points because they represent different trade-offs between computational efficiency and 3D geometric expressiveness \cite{yan2018second,shi2020pvrcnn,yin2021centerpoint}.

Detection performance will be evaluated using standard metrics such as precision, recall, and average precision. In addition, the dissertation will report center-position error, horizontal error, vertical error, size error, and detection performance as a function of distance and point count. These additional metrics are necessary because average precision alone does not show whether detections are accurate enough for tracking and trajectory prediction.

\section{Component 2: Tracking and State Estimation from Sparse LiDAR Detections}

\subsection{Purpose}

The second methodological component transforms per-frame detections into temporally stable object tracks. This is necessary because single-frame detections can be noisy, may disappear for short periods, and do not directly provide persistent object identity or velocity. A tracking module estimates the state of each drone over time and provides the input for trajectory prediction.

\subsection{Tracking-by-Detection Pipeline}

The tracking pipeline follows the tracking-by-detection paradigm. In each frame, the detector produces a set of 3D bounding boxes and confidence scores. The center of each predicted box is interpreted as a noisy measurement of a drone position. Existing tracks are predicted forward to the current timestamp, detections are assigned to tracks, and matched tracks are updated with the new measurements.

The first baseline will use nearest-neighbor association and a constant-velocity motion model. This provides a transparent reference implementation and helps identify the effect of detection noise and missed detections. The second baseline will use a Kalman filter for recursive state estimation \cite{kalman1960newApproach}. For each track, the state vector is defined as
\begin{math}
[
\mathbf{x}_t =
\begin{bmatrix}
p_x & p_y & p_z & v_x & v_y & v_z
\end{bmatrix}^{\top},
]
\end{math}

where ($p_x$, $p_y$, $p_z$) denote three-dimensional position and ($v_x$, $v_y$, $v_z$) denote velocity. A constant-velocity transition model predicts the next state, while the detector provides noisy position measurements. The Kalman filter then combines prediction and measurement correction to obtain a smoothed state estimate.

For multi-drone scenarios, detection-to-track association will first be implemented with distance-based nearest-neighbor matching. It will then be extended to an assignment formulation using the Hungarian algorithm \cite{kuhn1955hungarian}. This allows simultaneous association of multiple detections and tracks. The tracking design is inspired by efficient tracking-by-detection systems such as SORT and AB3DMOT, while adapting the state and measurement assumptions to sparse airborne LiDAR drone observations \cite{bewley2016sort,weng2020ab3dmot}.

\subsection{Uncertainty Handling}

A key methodological element is uncertainty handling. Sparse LiDAR detections of small drones can vary strongly in quality depending on range, viewing angle, and number of returns. Therefore, the tracking module will investigate whether measurement uncertainty can be adapted using detection confidence, number of LiDAR points inside the detection, or estimated spatial spread of the observed points.

In a Kalman-filter setting, this can be represented through the measurement covariance. A high-confidence detection with many supporting points can be treated as a more reliable measurement, while a low-confidence detection or a detection with very few points can be assigned higher measurement uncertainty. This allows the tracker to rely more strongly on the motion model when detections are weak, and more strongly on measurements when detections are reliable.

Missed detections will be handled by prediction-only updates for a limited number of frames. If no detection is associated with a track, the state will be propagated using the motion model and the uncertainty will increase. After a configurable number of missed frames, the track will be terminated. This mechanism is important because small drones may disappear temporarily due to sparse point returns, but tracks should not be kept indefinitely without evidence.

\subsection{Tracking Evaluation}

Tracking will be evaluated both independently and as input to prediction. The metrics will include:

\begin{itemize}
\item track recall and track precision,
\item number of identity switches,
\item average track duration,
\item position error over time,
\item velocity estimation error,
\item robustness to missed detections,
\item runtime per frame.
\end{itemize}

The evaluation will compare raw detector centers, nearest-neighbor tracking, Kalman-filter tracking, and assignment-based multi-object tracking. The purpose is to determine whether temporal filtering improves state estimation sufficiently to support trajectory prediction.

\section{Component 3: Short-Horizon Trajectory Prediction}

\subsection{Purpose}

The third methodological component predicts future drone positions from the estimated tracks. In a helicopter detect-and-avoid context, short-horizon prediction is necessary because avoidance decisions depend not only on where an intruder is now, but also on where it is likely to be in the near future. The dissertation therefore investigates how reliable future trajectory estimates can be produced from sparse LiDAR-derived tracks.

\subsection{Model-Based Prediction Baselines}

The first prediction methods will be simple model-based baselines. These include constant-position, constant-velocity, and Kalman-based prediction. A constant-velocity predictor estimates future position according to
\begin{math}
[
\mathbf{p}{t}
+
\mathbf{v}{t}\Delta t.
]
\end{math}

Here, ($\mathbf{p}{t}$) and ($\mathbf{v}{t}$) are provided by the tracking module. The Kalman-based variant additionally propagates the state covariance, which provides an estimate of prediction uncertainty. These baselines are computationally efficient, interpretable, and suitable for short prediction horizons. They also provide a necessary reference point for evaluating whether more complex learned models are justified.

\subsection{Learning-Based Prediction Models}

In a second step, learning-based predictors will be investigated. These models receive a sequence of past drone states and predict a sequence of future positions. The input may consist of positions only, or it may include velocities, detection confidence, point-count information, and track uncertainty. The first learning-based baselines will use simple sequence models trained on ground-truth trajectories. Later experiments will use tracks generated by the detection and tracking pipeline to evaluate robustness under realistic perception errors.

Existing work on non-cooperative UAV trajectory prediction shows that future motion can be predicted online from observed flight behavior \cite{xie2022noncooperativeUAVPrediction}. General multi-agent trajectory prediction methods such as Social-LSTM and Trajectron++ demonstrate how temporal history and interaction structure can be incorporated into learned forecasting models \cite{alahi2016socialLSTM,salzmann2020trajectronpp}. In this dissertation, such ideas will be adapted cautiously, with a focus on short-horizon prediction and sparse LiDAR-derived input rather than idealized ground-truth trajectories.

A central comparison will be made between predictors trained on clean ground-truth tracks and predictors evaluated on noisy detector-based tracks. This distinction is essential because a model that performs well on perfect trajectories may fail when the input contains missed detections, identity switches, or noisy height estimates.

\subsection{Prediction Uncertainty}

For safety-relevant perception, a single deterministic future trajectory may be insufficient. A drone may maneuver unexpectedly, and the perception system may be uncertain about its current state. Therefore, the methodology will include uncertainty-aware prediction. For Kalman-based methods, uncertainty is obtained through covariance propagation. For learning-based models, uncertainty may be estimated through probabilistic outputs, model ensembles, or sampling-based prediction.

The purpose of uncertainty estimation is not to solve the final avoidance problem, but to provide a richer input to a future conflict assessment module. A predicted position with high uncertainty should be interpreted differently from a predicted position with low uncertainty. This distinction is important for safety-oriented decision making.

\subsection{Trajectory Prediction Evaluation}

Prediction quality will be evaluated using standard forecasting metrics and application-specific metrics. The standard metrics include:

\begin{itemize}
\item average displacement error,
\item final displacement error,
\item horizontal displacement error,
\item vertical displacement error,
\item error as a function of prediction horizon.
\end{itemize}

In addition, the dissertation will evaluate how prediction quality depends on detection quality, track length, point sparsity, drone distance, and missed detections. For helicopter detect-and-avoid relevance, scenario-level quantities such as future separation and time-to-conflict-oriented indicators will be considered. These metrics will not replace a full avoidance system, but they will indicate whether the perception pipeline produces information that is useful for downstream safety logic.

\section{Experimental Procedure}

The experimental procedure will follow a staged design. First, synthetic LiDAR datasets will be generated and organized into training, validation, and benchmark splits. Second, baseline 3D detectors will be trained and evaluated. Third, tracking methods will be applied to detector outputs and compared against ground-truth trajectories. Fourth, trajectory prediction methods will be trained and evaluated using both ideal ground-truth tracks and realistic detector-based tracks. Finally, the complete pipeline will be evaluated on unseen benchmark scenarios.

The experiments will therefore separate four levels of performance:

\begin{enumerate}
\item single-frame detection performance,
\item tracking and state-estimation performance,
\item trajectory prediction performance from ground-truth tracks,
\item trajectory prediction performance from detector-based tracks.
\end{enumerate}

This separation is important because errors can propagate through the pipeline. A poor prediction may be caused by an unsuitable prediction model, but it may also be caused by detector noise, wrong data association, height errors, or missed detections. The staged evaluation allows these effects to be analyzed separately.

\section{Expected Contribution}

The expected contribution of the methodology is an integrated framework for sparse LiDAR-based drone tracking and trajectory prediction. The dissertation aims to contribute:

\begin{itemize}
\item a benchmark structure for small airborne drone encounters in LiDAR point clouds,
\item a systematic analysis of how sparse detections affect tracking and prediction,
\item a tracking pipeline for non-cooperative drones using LiDAR-based detections,
\item a comparison of model-based and learning-based short-horizon trajectory predictors,
\item an evaluation methodology that connects detection, tracking, and prediction performance to helicopter detect-and-avoid perception requirements.
\end{itemize}

The scientific value lies in bridging the gap between per-frame LiDAR object detection and safety-relevant temporal perception. Instead of treating drone detection as an isolated bounding-box prediction problem, the proposed methodology studies how sparse airborne LiDAR observations can support continuous intruder state estimation and future motion prediction.